%% ========================================================================
%% Chapter 5: Process Management
%% ========================================================================

\chapter{Process Management}
\label{ch:process-management}

\chapterhero{red}{Observe how processes are created, scheduled, prioritized, and stopped so you can control system workloads accurately.}{\faTasks\quad process}{\faBell\quad signal}{\faSortAmountDown\quad priority}

This chapter discusses the concepts of processes and threads in Linux, the
process lifecycle, scheduling, signals, job management, and real-time process
monitoring. Understanding process management is a core skill for Linux system
administrators to maintain server performance and stability.

\textbf{What You'll Learn}

By the end of this chapter, you will be able to:

\begin{enumerate}
    \item explain the concepts of processes and threads and the differences between them in Linux;
    \item read process status and understand the process lifecycle from creation to completion;
    \item adjust process priority using \cmd{nice} and \cmd{renice};
    \item send signals to processes using \cmd{kill}, \cmd{pkill}, and \cmd{killall};
    \item manage foreground and background jobs using \cmd{fg}, \cmd{bg}, \cmd{jobs}, and \cmd{nohup}.
\end{enumerate}

\begin{examplebox}[Think of Linux as a busy food stall]
Imagine a busy food stall. Some customers are eating, some are waiting for
their order, some are queuing at the counter, and some are having their orders
taken. Each customer is doing something different at roughly the same time:
that is exactly how Linux manages \textbf{processes}.

Every program running on Linux is like a customer: it has its own queue number
(\textbf{PID}), its own current activity (\textbf{status}), and it is served
in turns by the cashier (\textbf{CPU}). VIP customers get served faster
(\textbf{priority}), and if a customer causes trouble, the owner can remove
them without disturbing the other customers (\textbf{kill signal}).

There is no need to be intimidated by the technical terms : they are just
formal names for concepts that are already very familiar in everyday life.
\end{examplebox}

%% ------------------------------------------------------------------------
\section{Process and Thread Concepts}
\label{sec:process-concept}

\subsection{Processes}

Have you ever opened several apps on your smartphone at the same time : for
example, music keeps playing while you read a chat? That is a real-world
example of multiple \textbf{processes} running simultaneously. In Linux, the
concept is exactly the same.

When you run a program, Linux does not simply execute its instructions
directly. The operating system creates a \textbf{process} : an \textit{instance}
of the program being executed, complete with its own memory space, file
descriptors, and execution state. Every time a program is started, whether
from the terminal or automatically by the system, Linux creates a new process
with a unique PID (Process ID).

This isolation between processes is what makes Linux stable. If one process
crashes, other processes are unaffected because each has a completely separate
memory space. The analogy: if one browser tab crashes, the other tabs stay
open and keep running normally.

Every process has the following key attributes:

\begin{itemize}
  \item \textbf{PID}: Unique process identifier number.
  \item \textbf{PPID}: The PID of the parent process that created it.
  \item \textbf{UID/GID}: Owner identity, determines the process's access rights.
  \item \textbf{Status}: Current condition of the process (running, sleeping, stopped, etc.).
  \item \textbf{Nice}: Scheduling priority value (-20 to +19).
\end{itemize}

\subsection{Threads}

Within a single process, it is possible to have multiple execution flows
running concurrently : these are called \textbf{threads}. Unlike processes,
threads share memory space and resources with other threads in the same process.

Threads are useful when a program needs to perform multiple tasks in parallel,
for example a web server handling many requests simultaneously without having
to create a new process for each request. Creating a thread is far lighter than
creating a new process because no new memory space needs to be allocated.
However, sharing memory also means that one problematic thread can potentially
bring down the entire process.

\begin{examplebox}[Process and Thread Analogy]
  \begin{itemize}
    \item \textbf{Process} = a restaurant with its own kitchen. Two restaurants
      do not share ingredients : a problem in one restaurant does not affect
      the other.
    \item \textbf{Thread} = chefs working in one shared kitchen. They share
      the same stove and ingredients, making it more efficient, but if one
      chef spills hot oil, everyone in the kitchen is affected.
  \end{itemize}
\end{examplebox}

\subsection*{Lab 5.1 : Viewing Processes and Threads}

\begin{enumerate}
  \item Display all running processes:
\begin{lstlisting}[language=bash, caption={Displaying all processes}]
ps aux
\end{lstlisting}

  \item Display processes along with their threads, visible in the \texttt{LWP}
    (Light-Weight Process ID) column:
\begin{lstlisting}[language=bash, caption={Displaying processes with threads}]
ps aux -L
\end{lstlisting}

  \item View the PID of the active shell and its process details:
\begin{lstlisting}[language=bash, caption={Current shell process details}]
echo $$
ps -p $$ -f
\end{lstlisting}

  \item View the process hierarchy visually:
\begin{lstlisting}[language=bash, caption={Process tree}]
pstree -p
\end{lstlisting}
\end{enumerate}

\begin{exercisebox}[5.1]
  Run \cmd{ps aux} and observe its output:
  \begin{enumerate}
    \item How many processes are running in total? Which process has the
      smallest PID?
    \item Run \cmd{pstree -p} and find your \texttt{bash} process. What
      process is the parent (PPID) of that bash?
    \item Compare the output of \cmd{ps aux} and \cmd{ps aux -L}. What
      differences do you observe?
  \end{enumerate}
\end{exercisebox}

%% ------------------------------------------------------------------------
\section{Process Lifecycle}
\label{sec:process-lifecycle}

\subsection{Process States}

A process is not always being actively executed by the CPU throughout its
lifetime. Processes transition between several states depending on what is
happening : executing, waiting for I/O, paused, or finished. These states
can be seen in the \texttt{STAT} column of \cmd{ps} output.

Knowing the state of a process is very useful for diagnosing problems. A large
number of \texttt{D} processes indicates an I/O bottleneck, while an
accumulation of \texttt{Z} processes indicates a bug in the parent program
that is not cleaning up its child processes.

\begin{table}[htbp]
  \centering
  \caption{Process states in Linux}
  \begin{tabular}{lll}
    \hline
    \textbf{Code} & \textbf{Name} & \textbf{Description} \\
    \hline
    R & Running/Runnable & Being executed by the CPU or ready to execute \\
    S & Sleeping         & Waiting for an event; can be interrupted by a signal \\
    D & Uninterruptible  & Waiting for I/O; cannot be interrupted by a signal \\
    T & Stopped          & Paused by SIGSTOP signal or a debugger \\
    Z & Zombie           & Finished but not yet cleaned up by the parent process \\
    \hline
  \end{tabular}
\end{table}

\begin{notebox}
  A \textbf{zombie} process does not consume meaningful CPU or memory, but it
  still occupies a slot in the process table. If they accumulate in large
  numbers, the system can run out of available PIDs.
\end{notebox}

\subsection{The fork-exec Mechanism}

Every new process in Linux is born through two steps: \texttt{fork()} and
\texttt{exec()}. When you type a command in the shell, the shell calls
\texttt{fork()} : copying itself into a child process. This child process then
calls \texttt{exec()} : replacing its memory contents with the program to be
run.

This two-step approach is not arbitrary. By calling \texttt{fork()} first, the
child process inherits the complete configuration of the parent (environment
variables, open file descriptors) before the new program is launched. This is
what allows the shell to set up I/O redirection \textit{before} the program
executes.

\begin{tipbox}
  Check the exit code of the last command using the variable \var{\$?}. A
  value of \texttt{0} means success; any other value means there was an error.
  This is useful for writing scripts that respond to the success or failure of
  a command.
\begin{lstlisting}[language=bash]
ls /tmp && echo "Exit code: $?"      # prints 0
ls /does-not-exist ; echo "Exit code: $?" # prints 2
\end{lstlisting}
\end{tipbox}

\subsection*{Lab 5.2 : Observing the Process Lifecycle}

\begin{enumerate}
  \item Create a background process and observe its state:
\begin{lstlisting}[language=bash, caption={Observing process state}]
sleep 60 &
ps aux | grep sleep
\end{lstlisting}

  \item Observe the exit code of successful and failed commands:
\begin{lstlisting}[language=bash, caption={Observing exit codes}]
ls /tmp
echo "Success: $?"

ls /directory-does-not-exist
echo "Failed: $?"
\end{lstlisting}
\end{enumerate}

\begin{exercisebox}[5.2]
  \begin{enumerate}
    \item Run \cmd{sleep 120 \&} and observe the \texttt{STAT} column in
      \cmd{ps aux}. What state is shown? Why is the sleep process in that
      state?
    \item Run several successful and failing commands, then record the exit
      code of each. What pattern do you find?
  \end{enumerate}
\end{exercisebox}

%% ------------------------------------------------------------------------
\section{Process Scheduling and Priority}
\label{sec:process-scheduling}

\subsection{The Scheduler and Nice Values}

A single CPU core can only execute one process at a time, yet Linux is capable
of running hundreds of processes \textit{apparently simultaneously} thanks to
the \textbf{scheduler} : a kernel component that determines which process gets
CPU time next. Linux uses the CFS (Completely Fair Scheduler), which attempts
to distribute CPU time fairly among all processes.

The priority of each process is controlled by the \textbf{nice} value, ranging
from \textbf{-20} (highest priority) to \textbf{+19} (lowest priority), with a
default of \textbf{0}. Adjusting priority is very useful when you need to ensure
that important processes : such as a database service : always receive sufficient
CPU time, while preventing heavy processes like backups from degrading overall
system performance.

\begin{notebox}
  Only \textbf{root} can set a negative nice value (higher priority than
  default). A regular user can only increase the nice value (lower priority)
  for their own processes.
\end{notebox}

\subsection{The nice and renice Commands}

\cmd{nice} and \cmd{renice} are two sides of the same coin. \cmd{nice} is used
when \textit{starting} a new process with a specific priority value, while
\cmd{renice} is used to change the priority of a process that is \textit{already
running}. Use \cmd{nice} when you are about to start a heavy, non-urgent task
such as compressing a large file, and use \cmd{renice} when you realise a
running process needs its priority adjusted.

\begin{lstlisting}[language=bash, caption={Syntax of nice and renice}]
# Run a command with nice +10 (lower priority than normal)
nice -n 10 command

# Change the nice value of a running process
renice -n 15 -p <PID>
\end{lstlisting}

\subsection*{Lab 5.3 : Adjusting Process Priority}

\begin{enumerate}
  \item Run a process with low priority:
\begin{lstlisting}[language=bash, caption={Running a process with nice +10}]
nice -n 10 sleep 300 &
\end{lstlisting}

  \item Verify the nice value in the \texttt{NI} column:
\begin{lstlisting}[language=bash, caption={Viewing the nice value}]
ps aux | grep sleep
\end{lstlisting}

  \item Change the nice value of the running process:
\begin{lstlisting}[language=bash, caption={Changing nice with renice}]
renice -n 15 -p <PID>
ps -p <PID> -o pid,ni,cmd
\end{lstlisting}

  \item Clean up the test process:
\begin{lstlisting}[language=bash, caption={Stopping the test process}]
kill %1
\end{lstlisting}
\end{enumerate}

\begin{exercisebox}[5.3]
  \begin{enumerate}
    \item Run \cmd{nice -n 5 sleep 200 \&} and verify its NI value with
      \cmd{ps}.
    \item Change the nice value to 10 using \cmd{renice}, then verify again.
    \item Try to change the nice value to -5 without sudo. What happens? Why
      does Linux restrict this for regular users?
  \end{enumerate}
\end{exercisebox}

%% ------------------------------------------------------------------------
\section{Process Signals}
\label{sec:process-signals}

\subsection{Signals as a Communication Mechanism}

How do you communicate with a running process : for example to ask it to stop,
pause, or reload its configuration : without interacting directly with the
program? The answer is \textbf{signals}. A signal is an asynchronous
notification sent to a process to interrupt its execution and trigger an action.
Each signal has a name (e.g. SIGTERM) and a number (e.g. 15).

Signals are the only standard way to send instructions to a process from
outside the process itself : critical for processes running without an
interactive interface, such as background system services.

\begin{table}[htbp]
  \centering
  \caption{Important signals in Linux}
  \begin{tabular}{llll}
    \hline
    \textbf{Name} & \textbf{No.} & \textbf{Shortcut} & \textbf{Description} \\
    \hline
    SIGHUP  & 1  & :                & Reload configuration (hangup) \\
    SIGINT  & 2  & \keystroke{Ctrl+C} & Stop from keyboard \\
    SIGKILL & 9  & :                & Force stop; cannot be caught \\
    SIGTERM & 15 & :                & Request graceful stop (default) \\
    SIGSTOP & 19 & \keystroke{Ctrl+Z} & Pause process; cannot be caught \\
    SIGCONT & 18 & :                & Resume a paused process \\
    \hline
  \end{tabular}
\end{table}

\subsection{The kill, pkill, and pgrep Commands}

Despite its name, \cmd{kill} is not always used to terminate a process : its
primary function is to \textit{send a signal} to a process. Without a signal
argument, \cmd{kill} sends SIGTERM by default. \cmd{pkill} works similarly but
targets processes by name rather than PID. \cmd{pgrep} is useful for finding
PIDs by process name before sending a signal with \cmd{kill}.

\begin{lstlisting}[language=bash, caption={Using kill, pkill, and pgrep}]
# Send SIGTERM (default) to a specific PID
kill 1234

# Send SIGKILL to a specific PID
kill -9 1234

# Find PIDs by process name
pgrep sleep

# Send a signal to all processes matching a name
pkill sleep
pkill -9 process-name
\end{lstlisting}

\begin{tipbox}
  Always try \textbf{SIGTERM} first before \textbf{SIGKILL}. SIGTERM gives the
  process a chance to save data and clean up resources before stopping.
  SIGKILL forcefully terminates the process without any cleanup : risking data
  corruption or unreleased file locks.
\end{tipbox}

\subsection*{Lab 5.4 : Sending Signals to Processes}

\begin{enumerate}
  \item Create test processes:
\begin{lstlisting}[language=bash, caption={Creating test processes}]
sleep 500 &
sleep 600 &
sleep 700 &
ps aux | grep -v grep | grep sleep
\end{lstlisting}

  \item Stop one process with SIGTERM and verify:
\begin{lstlisting}[language=bash, caption={Stopping a process with SIGTERM}]
kill <PID-sleep-500>
ps aux | grep -v grep | grep sleep
\end{lstlisting}

  \item Pause and resume a process with SIGSTOP/SIGCONT:
\begin{lstlisting}[language=bash, caption={Pausing and resuming a process}]
kill -SIGSTOP <PID-sleep-600>
ps aux | grep sleep      # observe STAT column: changes to T

kill -SIGCONT <PID-sleep-600>
ps aux | grep sleep      # STAT returns to S
\end{lstlisting}

  \item Stop all sleep processes at once:
\begin{lstlisting}[language=bash, caption={Stopping all sleep processes}]
pkill sleep
\end{lstlisting}
\end{enumerate}

\begin{exercisebox}[5.4]
  \begin{enumerate}
    \item Run \cmd{sleep 400 \&}, send SIGSTOP, and observe the change in the
      \texttt{STAT} column. What state appears?
    \item Send SIGCONT and verify the process is running again.
    \item Stop the process with SIGTERM then verify it is gone. When would
      you choose SIGKILL over SIGTERM?
  \end{enumerate}
\end{exercisebox}

%% ------------------------------------------------------------------------
\section{Job Management}
\label{sec:job-management}

\subsection{Foreground and Background}

In the context of a shell, a \textbf{job} is a process managed by that shell.
Each job can run in the \textbf{foreground} : holding the terminal so you
cannot type other commands : or in the \textbf{background} : running behind the
scenes while the terminal remains free to use.

Append \texttt{\&} to the end of a command to run it in the background. This
is very useful for long-running tasks such as file compression or system
updates, so you do not have to wait and can continue other work. Without this
mechanism, every long-running task would block the terminal and force you to
open a new session just to run the next command.

\begin{table}[htbp]
  \centering
  \caption{Job management commands and shortcuts}
  \begin{tabular}{ll}
    \hline
    \textbf{Command/Shortcut} & \textbf{Function} \\
    \hline
    \texttt{command \&}    & Run command in the background \\
    \cmd{jobs}             & List active jobs in this shell \\
    \cmd{fg \%N}           & Bring job number N to the foreground \\
    \cmd{bg \%N}           & Resume job number N in the background \\
    \keystroke{Ctrl+Z}     & Pause the currently running foreground job \\
    \keystroke{Ctrl+C}     & Stop the currently running foreground job \\
    \cmd{kill \%N}         & Stop job number N \\
    \hline
  \end{tabular}
\end{table}

\begin{notebox}
  Job references using \texttt{\%N} (job number within this shell) are
  different from PID (system-wide unique number). Both can be used with
  \cmd{kill}: \cmd{kill \%1} or \cmd{kill 4521} refer to the same thing if
  job 1 has PID 4521.
\end{notebox}

\subsection*{Lab 5.5 : Foreground and Background Job Management}

\begin{enumerate}
  \item Start three jobs in the background:
\begin{lstlisting}[language=bash, caption={Creating three background jobs}]
sleep 200 &
sleep 300 &
sleep 400 &
jobs
\end{lstlisting}

  \item Bring the first job to the foreground, pause it, then send it back
    to the background:
\begin{lstlisting}[language=bash, caption={Moving a job between foreground and background}]
fg %1
# Press Ctrl+Z to pause
bg %1
jobs
\end{lstlisting}

  \item Stop all jobs:
\begin{lstlisting}[language=bash, caption={Stopping all jobs}]
kill %1 %2 %3
jobs
\end{lstlisting}
\end{enumerate}

\begin{exercisebox}[5.5]
  \begin{enumerate}
    \item Run \cmd{top} in the foreground. What happens in the terminal?
    \item Press \keystroke{Ctrl+Z} and check its status with \cmd{jobs}.
      What state is shown?
    \item Move it to the background with \cmd{bg}. Can \cmd{top} run
      properly in the background? Why?
    \item Bring it back to the foreground with \cmd{fg}, then exit with
      \keystroke{q}.
  \end{enumerate}
\end{exercisebox}

%% ------------------------------------------------------------------------
\section{Process Monitoring}
\label{sec:process-monitoring}

\subsection{The ps Command}

\cmd{ps} displays a \textit{snapshot} of process states at the moment the
command is run. Because its output is static, \cmd{ps} is well suited for
scripting : its output can be further processed with \cmd{grep}, \cmd{awk},
or saved to a file for later analysis.

\begin{lstlisting}[language=bash, caption={Variations of ps usage}]
# All processes with full format (most common)
ps aux

# Display process hierarchy visually
ps aux --forest

# Sort by CPU or memory usage
ps aux --sort=-%cpu | head -10
ps aux --sort=-%mem | head -10

# Specific information from a single process
ps -p 1234 -o pid,ppid,ni,stat,cmd
\end{lstlisting}

\subsection{The top Command}

While \cmd{ps} only takes a snapshot, \cmd{top} monitors processes in
\textit{real-time} with a display that refreshes every few seconds. Use
\cmd{top} when you need to monitor the system live, especially to find
processes that are suddenly consuming excessive CPU or memory : something a
\cmd{ps} snapshot would not capture.

Useful keyboard shortcuts inside \cmd{top}:

\begin{itemize}
  \item \keystroke{P}: Sort by CPU. \quad \keystroke{M}: Sort by memory.
  \item \keystroke{k}: Send a signal to a process. \quad \keystroke{r}: Change nice value.
  \item \keystroke{u}: Filter by user. \quad \keystroke{1}: Show each CPU core.
  \item \keystroke{q}: Quit.
\end{itemize}

\subsection{The htop Command}

\cmd{htop} is a modern version of \cmd{top} with a coloured display, CPU and
memory bar graphs, and interactive navigation using both keyboard and mouse.
The bar graphs make it easy to spot bottlenecks at a glance without having to
interpret numbers : just look at which bar is full. Press \keystroke{F9} to
send a signal directly to the selected process without needing to look up its
PID first.

\begin{lstlisting}[language=bash, caption={Installing and using htop}]
sudo apt install htop
htop
\end{lstlisting}

\subsection*{Lab 5.6 : Process Monitoring}

\begin{enumerate}
  \item Find the processes with the highest CPU and memory usage:
\begin{lstlisting}[language=bash, caption={Processes with the highest resource usage}]
ps aux --sort=-%cpu | head -10
ps aux --sort=-%mem | head -10
\end{lstlisting}

  \item Run \cmd{top} and explore its shortcuts:
\begin{lstlisting}[language=bash, caption={Exploring top}]
top
# Press M, P, 1, u in turn
# Press q to quit
\end{lstlisting}

  \item Install and run \cmd{htop}:
\begin{lstlisting}[language=bash, caption={Using htop}]
sudo apt install -y htop
htop
# Press F6 to select the sort column
# Press F10 or q to quit
\end{lstlisting}
\end{enumerate}

\begin{exercisebox}[5.6]
  \begin{enumerate}
    \item Use \cmd{ps aux --sort=\%mem} to find the process using the most
      memory on your VM. What process is it?
    \item Inside \cmd{top}, press \keystroke{1}. What changes in the display?
      Why is this information useful?
    \item Inside \cmd{htop}, navigate to the \texttt{sshd} process using the
      arrow keys. Press \keystroke{F9} and observe the signal options
      available.
  \end{enumerate}
\end{exercisebox}

%% ------------------------------------------------------------------------
\section{Summary}
\label{sec:process-summary}

\begin{itemize}
  \item A \textbf{process} is an instance of a running program with a unique
    PID and a separate memory space. \textbf{Threads} share memory within one
    process : lighter weight but more susceptible to mutual interference.
  \item \textbf{Process lifecycle}: R (running), S (sleeping), D (uninterruptible),
    T (stopped), Z (zombie). Processes are created via \texttt{fork-exec} and
    end with an exit code (0 = success).
  \item \textbf{Priority} is controlled by the nice value (-20 to +19, default 0).
    \cmd{nice} for new processes, \cmd{renice} for already-running processes.
    Negative values can only be set by root.
  \item \textbf{Signals} are the way to communicate with running processes.
    Always try SIGTERM (15) before SIGKILL (9). Use \cmd{kill} (by PID),
    \cmd{pkill} (by name), or \cmd{pgrep} to find PIDs.
  \item \textbf{Job management}: append \texttt{\&} for background, \cmd{fg}/\cmd{bg}
    to switch, \keystroke{Ctrl+Z} to pause.
  \item \textbf{Monitoring}: \cmd{ps} for snapshots and scripting, \cmd{top} for
    real-time monitoring, \cmd{htop} for a more visual interactive interface.
\end{itemize}

%% ------------------------------------------------------------------------
\section{Exercises}
\label{sec:process-exercises}

\begin{exercisebox}[5.A]
  \textbf{Exploring System Processes}
  \begin{enumerate}
    \item Run \cmd{ps aux --forest} and find the process with PID 1. What is
      the name and function of this process in a modern Linux system?
    \item Count how many processes are owned by user \texttt{root} and how
      many are owned by your user. Why does root have more processes?
    \item Find all processes in the \texttt{S} state. Why do most processes
      on a system reside in this state?
  \end{enumerate}
\end{exercisebox}

\begin{exercisebox}[5.B]
  \textbf{Job Management Simulation}
  \begin{enumerate}
    \item Run three \cmd{sleep} commands with durations of 100, 200, and 300
      seconds in the background. Verify all three with \cmd{jobs}.
    \item Bring the second job to the foreground, pause it with
      \keystroke{Ctrl+Z}, then send it back to the background with \cmd{bg}.
    \item Stop the first job with \cmd{kill \%1}. Display the job list again.
      How many jobs remain?
  \end{enumerate}
\end{exercisebox}

\begin{exercisebox}[5.C]
  \textbf{Priority and Signals}
  \begin{enumerate}
    \item Run two sleep processes: one with nice +5 and one with nice +15.
      Verify the NI value of both with \cmd{ps}.
    \item Use \cmd{renice} to change the first process's nice to +10. Which
      process does the scheduler now prioritise?
    \item Send SIGSTOP to one of the processes, verify its \texttt{T} state,
      then send SIGCONT. Terminate all test processes with \cmd{pkill sleep}.
  \end{enumerate}
\end{exercisebox}

%
